\documentclass[12pt]{article}
\usepackage[utf8]{inputenc}
\usepackage[T1]{fontenc}
\usepackage[margin=1in]{geometry}
\usepackage{setspace}
\usepackage{xcolor}
\usepackage[colorlinks=true, urlcolor=blue!60!black, linkcolor=blue!60!black]{hyperref}
\setlength{\parindent}{0em}
\setlength{\parskip}{0.45em}
\renewcommand{\baselinestretch}{1.0}
\begin{document}
\begin{center}
{\Large \textit{NLRB v. Jones \& Laughlin Steel Corp.}, 301 U.S. 1 (1937)\par}
\vspace{0.3em}
{\normalsize Case Brief --- One-Page Edition\par}
\end{center}

\textbf{Memory Jogger.} The 1937 Commerce Clause turn: the Wagner Act reaches an integrated steel manufacturer (\href{https://supreme.justia.com/cases/federal/us/301/1/}{301 U.S. 1}).

\textbf{Facts.} Jones \& Laughlin, the fourth-largest steel producer, ran mines, ships, railroads, and plants in many states and shipped about 75\% of its steel interstate. It fired ten union leaders at Aliquippa, and the NLRB ordered reinstatement with back pay.

\textbf{Procedural History.} The Fifth Circuit refused enforcement; the Supreme Court reversed.

\textbf{Votes.} 5--4. Hughes wrote for the Court; McReynolds dissented for the ``Four Horsemen.''

\textbf{Holding.} Congress may regulate intrastate activities, including manufacturing labor relations, closely and substantially related to interstate commerce; the NLRA is valid as applied.

\textbf{Reasoning.} ``It is the effect upon commerce, not the source of the injury, which is the criterion.'' A stoppage at so far-flung an enterprise would be immediate and might be catastrophic; \textit{Schechter} and \textit{Carter Coal} were distinguished. Hughes warned that federal power still may not erase the line between national and local. McReynolds insisted manufacturing is not commerce.

\textbf{Examples.} \emph{Same-Side:} a garment maker buying and selling interstate. \emph{Opposite-Side:} a ban on guns near schools (\textit{Lopez}). \emph{Fence-Sitter:} a neighborhood bakery selling only locally.

\textbf{Critique.} Originalists say the effects test erased enumerated powers; progressives praise its economic realism, and historians debate the ``switch in time.''

\textbf{Key Quotation.} ``Although activities may be intrastate in character when separately considered, if they have such a close and substantial relation to interstate commerce that their control is essential or appropriate to protect that commerce from burdens and obstructions, Congress cannot be denied the power to exercise that control.'' --- Chief Justice Hughes, Opinion of the Court, 301 U.S. at 37.

\end{document}
